\chapter{绪论}
在原有的 zjuthesis 模板~\cite{zjuthesis}基础上开发而来。
同学们使用时要注意对照模板与要求，切不可盲目使用。

\section{研究背景及意义}

如果你在Overleaf上编译本模板，请注意如下事项：

\begin{itemize}
    \item 删除根目录的 ``.latexmkrc'' 文件，否则编译失败且不报任何错误
    \item 字体有版权所以本模板不能附带字体，请务必手动上传字体文件，并在各个专业模板下手动指定字体。
        
\end{itemize}


\section{研究现状}

我们可以用includegraphics来插入现有的jpg等格式的图片，
如\autoref{fig:zju-logo}所示。

\begin{figure}[htbp]
    \centering
    \includegraphics[width=.3\linewidth]{logo/zju}
    \caption{\label{fig:zju-logo}大学LOGO}

\end{figure}

\par 如\autoref{tab:sample}所示，这是一张自动调节列宽的表格。

\begin{table}[htbp]
    \caption{\label{tab:sample}自动调节列宽的表格}
    \begin{tabularx}{\linewidth}{c|X<{\centering}}
        \hline
        第一列 & 第二列 \\ \hline
        xxx & xxx \\ \hline
        xxx & xxx \\ \hline
        xxx & xxx \\ \hline
    \end{tabularx}
\end{table}


\par 如\autoref{equ:sample}，这是一个公式

\begin{equation}
    \label{equ:sample}
    A=\overbrace{(a+b+c)+\underbrace{i(d+e+f)}_{\text{虚数}}}^{\text{复数}}
\end{equation}

\section{研究内容}

\section{论文结构}

\chapter{相关技术}

\begin{figure}[htbp]
    \centering
    \includegraphics[width=.3\linewidth]{example-image-a}
    \caption{\label{fig:fig-placeholder}图片占位符}
\end{figure}
